\documentclass[11pt]{article}
\usepackage[margin=1in]{geometry}
\usepackage{amsmath}
\usepackage{amssymb}
\usepackage{bm}
\usepackage{xurl}
\usepackage[colorlinks=true,linkcolor=blue,citecolor=blue,urlcolor=blue]{hyperref}

\title{Ghana Weather Anomalies and the ICCO Cocoa Price Indicator During the 2023--2024 El Ni\~no Episode}
\author{Siyan Dong}
\date{\today}

\begin{document}
\maketitle

\begin{abstract}
This note records an empirical design for using Ghana weather anomalies and the ICCO cocoa daily price indicator as a motivating application for the sieve-orthogonalized empirical likelihood framework. The purpose is not to report cleaned empirical results. Instead, the goal is to show why the setting is economically desirable, how the variables map into the formal model, and what data-audit and diagnostic checks are needed before reporting estimates.
\end{abstract}

\section{Purpose}

The formal paper studies the predictive-regression model
\begin{equation}
    Y_t=m(W_{t-1})+\beta X_{t-1}+U_t,
    \label{eq:design-model}
\end{equation}
where \(X_{t-1}\) is a persistent market predictor and \(W_{t-1}\in\mathbb R^d\) is a stationary exogenous covariate vector entering through an unknown smooth function \(m\). A useful empirical design is cocoa-market prediction during the 2023--2024 El Ni\~no episode. In this application, Ghana weather anomalies provide predetermined physical supply-risk variables, while the ICCO cocoa daily price indicator provides a global cocoa-market price benchmark.

The design should be presented as a motivating application, not as the core mathematical setting. The theorem remains generic. The empirical discussion only explains why the assumptions are plausible and what diagnostics should be checked.

\section{Empirical Mapping}

Let \(P_t^{ICCO}\) denote the ICCO cocoa daily price indicator in U.S. dollars per tonne. The return object is
\begin{equation}
    Y_t=\Delta\log P_t^{ICCO}
        =\log P_t^{ICCO}-\log P_{t-1}^{ICCO}.
\end{equation}
The persistent predictor is a lagged ICCO price-state variable, for example
\begin{equation}
    X_{t-1}=\log P_{t-1}^{ICCO},
\end{equation}
or a bounded transformation \(w(X_{t-1})\) used in the empirical-likelihood score. This choice keeps the application close to the local data in \texttt{data/processed/cocoa\_ghana\_full.csv}, which already contains \texttt{log\_price}, \texttt{log\_price\_lagt}, and \texttt{log\_return}.

The weather vector should be kept low-dimensional:
\begin{equation}
    W_{t-1}=
    \begin{pmatrix}
        A^{rain}_{t-1}\\
        A^{temp}_{t-1}
    \end{pmatrix}
    \in\mathbb R^2,
    \label{eq:weather-vector}
\end{equation}
where \(A^{rain}_{t-1}\) is a Ghana rainfall anomaly and \(A^{temp}_{t-1}\) is a Ghana average-temperature anomaly. In the processed file these correspond to \texttt{PRCP\_anom\_mean} and \texttt{TAVG\_anom\_mean}. Keeping \(d=2\) is important because the sieve rate window
\[
    T^{d/(2s)}\ll K\ll T^{1/2}
\]
is easier to defend when the nuisance dimension is small. Adding cross-station dispersion measures may be useful in robustness checks, but it raises the effective dimension of \(W\).

\section{Why Weather Anomalies}

Raw temperature and rainfall series contain strong seasonal patterns. Those seasonal components are partly anticipated by farmers, traders, and inventory managers, so they are not clean shocks. Weather anomalies are preferable because they remove the normal seasonal climatology and focus the design on abnormal heat or rainfall realizations.

This matters for the formal assumptions. The paper does not need raw weather to be independent of cocoa prices. It needs the projected sieve component of the centered market weight to be small. Anomaly construction makes this more plausible for two reasons. First, subtracting the climatological seasonal component reduces deterministic seasonality, making stationarity a more realistic approximation. Second, local weather anomalies are generated by physical atmospheric processes and are predetermined relative to next-period price changes, so they are more credible as exogenous controls than raw weather levels.

This is still an empirical claim to diagnose, not a fact to assert. In the note and paper, the correct language is:
\[
    \mathbb E\left[
        \bm R_K(W_{t-1})
        \{w(X_{t-1})-\mathbb E w(X_{t-1})\}
    \right]\approx \bm0,
\]
not ``\(W_t\) and \(X_t\) are independent.'' If the sample projection of the ICCO weight onto the weather-anomaly sieve is large, the empirical design should report that leakage and either modify the control set or treat the application as only suggestive.

\section{Economic Desirability}

The setting is economically desirable because the supply channel is direct. Ghana is a major cocoa producer, and abnormal heat or rainfall conditions in Ghana can affect expected cocoa supply through pod development, disease pressure, drying conditions, and harvest timing. The response is unlikely to be linear: moderate deviations may have limited effects, while extreme heat or rainfall shortages can have disproportionate consequences. This is exactly the kind of relationship represented by the unknown function \(m(W_{t-1})\).

The ICCO daily price indicator is a useful market object because it summarizes international cocoa-price conditions without requiring a futures-contract roll rule. The empirical question is then economically sharp:
\[
    \begin{gathered}
    \text{Does the lagged ICCO price state predict future ICCO returns}\\
    \text{after flexibly absorbing Ghana weather-anomaly effects?}
    \end{gathered}
\]
If \(\beta\neq0\), persistent market-state information contains incremental predictive content beyond the nonlinear weather-supply component. If \(\beta=0\), the apparent price-state signal may be explained by the weather-anomaly state and noise.

For an advisor or referee, this application also explains why a fixed-intercept predictive regression is too restrictive. During an event such as the 2023--2024 El Ni\~no episode, the intercept-like component of expected returns can move with nonlinear weather stress. Profiling out \(m(W_{t-1})\) gives the financial predictor a cleaner interpretation.

\section{Data Inventory and Audit}

The current local reference file is
\[
    \texttt{data/processed/cocoa\_ghana\_full.csv}.
\]
It merges Ghana weather anomalies with the ICCO price indicator and contains the fields needed for the design:
\[
\begin{array}{ll}
\texttt{date} & \text{observation date},\\
\texttt{PRCP\_anom\_mean} & \text{mean Ghana rainfall anomaly},\\
\texttt{TAVG\_anom\_mean} & \text{mean Ghana average-temperature anomaly},\\
\texttt{N\_stations} & \text{number of reporting weather stations},\\
\texttt{log\_price} & \log P_t^{ICCO},\\
\texttt{log\_price\_lagt} & \log P_{t-1}^{ICCO},\\
\texttt{log\_return} & \Delta\log P_t^{ICCO}.
\end{array}
\]
The June 2023--June 2024 event window is the preferred baseline window for discussion. The merged full file extends beyond that window, so the post-event observations can be used later for sensitivity checks or short-horizon aftermath analysis.

\paragraph{Data audit required.}
The raw ICCO file \texttt{data/raw/Daily Prices\_ICCO.csv} contains duplicate dates. Some duplicates are exact repeats, such as 15 December 2023 and 9 January 2024. More importantly, 30 January 2024 and 31 January 2024 contain conflicting price values. Because these conflicts mechanically generate implausible daily log returns in the processed file, no return estimates, event magnitudes, or regression results should be reported until the duplicate ICCO quotes are validated against the source. The cleaning rule should be documented before the empirical application is promoted from a design note to a results section.

\section{Diagnostics Before Estimation}

Before reporting Sieve-EL results, the empirical section should include four checks.

\paragraph{One price per date.}
Validate the ICCO price series, remove exact duplicates, and resolve conflicting duplicate quotes using the official source. Recompute \(\log P_t^{ICCO}\), \(Y_t\), and \(X_{t-1}\) only after this audit.

\paragraph{Anomaly behavior.}
Plot the rainfall and temperature anomalies over the full sample and the June 2023--June 2024 window. Check that the anomaly series are centered and that their autocorrelations decay. Unit-root tests can be reported as descriptive diagnostics, but the text should avoid claiming that finite-sample tests prove alpha mixing.

\paragraph{Projection leakage.}
For each candidate sieve dimension \(K\), project the centered ICCO weight \(w(X_{t-1})-\bar w_T\) onto the weather-anomaly sieve. Report the projection \(R^2\) and the norm \(\|T^{-1}\bm R'\widetilde{\bm w}\|\). The design is most convincing when this projection is small relative to the \(\sqrt{K/T}\) benchmark used in the theory.

\paragraph{Economic event plot.}
After the ICCO audit, plot Ghana rainfall and temperature anomalies together with the ICCO log price or return series, shading June 2023--June 2024. This plot is descriptive only. Its purpose is to show economic relevance and timing, not to establish causality.

\section{How This Fits the Paper}

The main paper should keep the theorem generic. The Ghana-ICCO application can be used in the introduction or interpretation section to explain the empirical need for a nonlinear exogenous-control function. The correct bridge is:
\[
    \text{weather anomaly control}
    \quad\longrightarrow\quad
    m(W_{t-1})
    \quad\longrightarrow\quad
    \text{sieve profiling and orthogonalized EL}.
\]
The application should not be used to claim full independence between weather and market predictors. Instead, it motivates the weaker and more testable orthogonality condition. This is more credible empirically and more consistent with the formal proof.

\section{Conclusion}

The Ghana-ICCO setting is a strong motivating example because it combines a visible commodity-market episode, a direct supply-risk channel, and a persistent price-state predictor. Weather anomalies make the exogenous-control argument more credible than raw weather levels. The low-dimensional vector in \eqref{eq:weather-vector} also keeps the sieve condition interpretable. The immediate next step is not estimation; it is to audit the ICCO duplicate dates and then run stationarity, projection-leakage, and event-window diagnostics.

\begin{thebibliography}{9}

\bibitem{icco-statistics}
International Cocoa Organization.
\newblock Statistics: cocoa prices and market data.
\newblock \url{https://www.icco.org/statistics/}.

\bibitem{ncei-ghcnd}
NOAA National Centers for Environmental Information.
\newblock Global Historical Climatology Network Daily.
\newblock \url{https://www.ncei.noaa.gov/products/land-based-station/global-historical-climatology-network-daily}.

\bibitem{chirps}
Climate Hazards Center, UC Santa Barbara.
\newblock CHIRPS: Rainfall Estimates from Rain Gauge and Satellite Observations.
\newblock \url{https://www.chc.ucsb.edu/data/chirps}.

\bibitem{wwa-west-africa-heat}
World Weather Attribution.
\newblock Dangerous humid heat in southern West Africa about \(4^\circ\)C hotter due to climate change.
\newblock \url{https://www.worldweatherattribution.org/dangerous-humid-heat-in-southern-west-africa-about-4c-hotter-due-to-climate-change/}.

\bibitem{noaa-enso-june-2023}
NOAA Climate.gov.
\newblock June 2023 ENSO update: El Ni\~no is here.
\newblock \url{https://www.climate.gov/news-features/blogs/enso/june-2023-enso-update-el-nino-here}.

\end{thebibliography}

\end{document}
